\documentclass[UTF8,12pt,a4paper]{ctexart}

\usepackage{amsmath,amssymb}
\usepackage{booktabs}
\usepackage{geometry}
\usepackage{graphicx}
\usepackage[hidelinks]{hyperref}
\geometry{left=2.25cm,right=2.25cm,top=3cm,bottom=1.85cm}
\ctexset{section={format=\centering\heiti\zihao{4}},subsection={format=\heiti\zihao{-4}}}
\pagestyle{plain}
\linespread{1}
\newcommand{\keywords}[1]{\par\noindent\textbf{关键词：}#1}

% Non-official LaTeX adaptation for a standalone Q3 chapter.
% Page dimensions follow the converted 2026 official Word template.
\title{\heiti\zihao{3} 算力约束下质量与规模的联合配置及结构转移分析}
\author{}
\date{}


\usepackage{tabularx,array,float,placeins}
\usepackage[font=normalsize]{caption}
\ctexset{section={number=\arabic{section}},subsection={number=\arabic{section}.\arabic{subsection}}}
\renewenvironment{abstract}{\begin{center}\heiti\zihao{3}摘\quad 要\end{center}\zihao{-4}}{\par}
\begin{document}
\maketitle
\begin{abstract}
针对有限算力下参数量、数据量、领域配比与治理质量的配置问题，考虑质量处理对训练Token的挤占及长上下文的额外开销。承接前两问的17域质量与13域响应，以实际训练配方凸包限制配比搜索，将质量收益作为需校准的桥接情景，建立三项成本一致的联合优化模型。

从预算可行方向推导质量投入判据，并得到规模弹性比等于一减治理成本份额的局部必要关系。采用多起点约束求解、真实配方锚点扩展及独立求解器对照，以四组消融区分规模、配比、治理与联合调整；同时用近等价搜索、预算加密和情景一致性排查表面结构变化。

原同域联合配置的预测平均Loss相对内部规模基线改善4.16\%至5.81\%，但165个场景中162个至少一个评测域退步。注意力与训练成本相当的临界上下文长度为30000。对齐质量提升端点价格后，55组中19组的最优成本形式改变；质量上界处治理份额下降也不代表绝对投入减少。按70\%情景支持规则未发现稳健联合结构转移。

结果提供三个预算量级的N/D/Q及完整17域配方，并区分有限候选推荐、机制比较与稳健情景选择。结论限于交接响应与治理假设；半合成质量数据、规模外推及上游独立验证资料不足，使其不能直接解释为真实训练收益或全局最优保证。
\keywords{算力约束；质量治理；资源替代；领域配比；结构性转移}
\end{abstract}
\clearpage
\section{问题分析、数据与基本假设}

问题三要求在算力约束下同时考虑参数规模、训练数据量、领域配比与质量治理，并判断预算跨越多个数量级时策略是否发生结构性变化。规模与数据的算力分配已有经验研究基础\cite{ref1}，但题给质量处理和注意力开销会改变资源替代关系。因此，不能仅沿用训练成本约束，也不能预设高预算必然对应更高质量投入份额。本文承接前两问的可执行响应与质量数据，先建立条件优化模型，再分开检验机制收益、数值求解和结论稳健性。总体关系见图~\ref{fig:q3_1}。


\begin{figure}[htbp]
\centering
\includegraphics[width=.95\textwidth]{figs/flow_overall_model.pdf}
\caption{质量、配方及规模响应的衔接与联合求解流程}\label{fig:q3_1}
\end{figure}

基本假设为：前两问交付参数在本次条件分析中固定；13个评测域等权，不预设业务优先级；治理质量的收益通过跨数据源桥接表达，该桥接尚未独立校准；训练、注意力和质量处理使用相同Token量；上下文长度由架构与任务外生确定；配比仅在实际训练配方的凸包内变化。N与D的上下界是公开的计算范围，而非附件识别出的自然极限，其影响另做敏感性检验。主要符号见表~\ref{tab:q3_1}。


\begin{table}[htbp]\centering\setlength{\tabcolsep}{3pt}\renewcommand{\arraystretch}{1.12}
\caption{主要符号与单位}\label{tab:q3_1}
\begin{tabularx}{\textwidth}{>{\raggedright\arraybackslash}p{2.975132cm}>{\raggedright\arraybackslash}p{8.925397cm}>{\raggedright\arraybackslash}X}\toprule
符号 & 含义 & 单位 \\
\midrule
$N, D$ & 参数个数、训练Token总数 & 个、Token \\
$n, d$ & N与D分别除以十亿后的数值 & 十亿参数、十亿Token \\
$p_i, q_i$ & 第i训练域配比、交接的固定基础质量 & 无量纲 \\
$q_{\mathrm{mix}}$ & 按配比加权的基础质量诊断量 & 无量纲 \\
$Q_0, Q$ & 治理基线质量、治理情景中的决策质量 & 无量纲 \\
$L_k, J$ & 第k评测域Loss、13域平均Loss & nats \\
$C, C_{\mathrm{used}}$ & 给定预算、实际计算成本 & FLOPs \\
$L_{\mathrm{ctx}}, \eta$ & 外生上下文长度、注意力成本系数 & Token、题给系数 \\
$a, h(Q)$ & 训练与注意力合并系数、单位Token增量治理成本 & 系数、FLOPs/Token \\
$s_Q, \kappa$ & 治理占实际成本比例、质量收益桥接强度 & 无量纲 \\
$\tau, \epsilon$ & 有效配比支持阈值、近等价Loss容差 & 无量纲、nats \\
\bottomrule\end{tabularx}
\end{table}

数据衔接见表~\ref{tab:q3_2}。训练领域17项质量实际只有6个不同取值，反映跨体系映射后的信息粒度；不能把这些值解释成17个独立估计的治理潜力。定义以下三个量，并将治理决策Q与配比加权诊断量区分：


\begin{equation}\label{eq:q3_1}
n=N/10^9,\quad d=D/10^9,\quad q_{\mathrm{mix}}=\sum_{i=1}^{17}p_iq_i,\quad Q_0=\frac{1}{17}\sum_{i=1}^{17}q_i.
\end{equation}

主情景取$Q_0=0.6219326581278445$；训练平均配方加权的0.6193580511629503作为另一基线敏感性情景。Q改变的是假设的总体治理水平，本文不直接修改固定领域质量向量。各领域的可达提升幅度尚无观测，因而不建立未经识别的逐域治理函数。


\begin{table}[htbp]\centering\setlength{\tabcolsep}{3pt}\renewcommand{\arraystretch}{1.12}
\caption{输入信息及其解释边界}\label{tab:q3_2}
\begin{tabularx}{\textwidth}{>{\raggedright\arraybackslash}p{2.578448cm}>{\raggedright\arraybackslash}p{7.437831cm}>{\raggedright\arraybackslash}X}\toprule
来源 & 实际使用的信息 & 解释边界 \\
\midrule
第一问交接 & 17域质量，约0.4785至0.7095；A4的512行训练配方 & 质量仅6种取值；凸包保证配比几何支持 \\
第二问交接 & 34维配方特征、13域响应、规模参数及求值器 & 沿用参数；原拟合脚本和训练/留出ID未完整交付 \\
B7质量部分 & 450条半合成记录所给质量标度律参数 & 跨数据源收益桥接未独立校准 \\
C7 & 2048、4096、8192、32768、131072 & 仅作为外生上下文的可行取值 \\
规模支持 & 1M/1B、60M/1B、1B/25B三组N/D尺度 & 本次正式优化点均按规模条件外推解释 \\
\bottomrule\end{tabularx}
\end{table}

配方回归及跨规模迁移的研究为接口组织提供方法背景\cite{ref2}，不替代本题实测证据。图~\ref{fig:q3_2}同时展示直接支持尺度与优化配置：处于512行配方凸包内，仍不代表任意N/D处的预测都获观测支持。A6至A11仅用于交接函数残差复算；在上游划分未知时，不能称为本问的新独立留出验证。


\begin{figure}[htbp]
\centering
\includegraphics[width=.95\textwidth]{figs/raw_q3_scale_support.pdf}
\caption{三个直接支持尺度与原联合配置的规模位置；所有优化结果按条件外推解释}\label{fig:q3_2}
\end{figure}

\section{质量感知的联合资源模型}

\subsection{配方与规模响应}

令i对应17个训练域、k对应13个评测域。对配比加入0.001平滑量后构造成分对数比，再与固定质量加权配比组成34维特征；使用第二问交付的均值、尺度和回归参数，不重新拟合：


\begin{equation}\label{eq:q3_2}
\operatorname{clr}_{\delta}(p)_i=\ln(p_i+\delta)-\frac{1}{17}\sum_{j=1}^{17}\ln(p_j+\delta),\quad\delta=0.001.
\end{equation}

\begin{equation}\label{eq:q3_3}
x=(p\odot q,\operatorname{clr}_{\delta}(p)),\quad z=(x-\mu)/\sigma,\quad\Phi_k=s_k(\beta_{0k}+\theta_k^\top z-c_k).
\end{equation}

特征按前17个质量加权配比、后17个成分对数比拼接，乘积和标准化按分量进行。式中回归截距、系数、均值、尺度、响应中心与响应斜率均取交接参数；零中心化的配方项表示相对响应。规模幅度与总响应为：


\begin{equation}\label{eq:q3_4}
S_k=(n/0.001)^{-\eta_k}(d/1)^{-\zeta_k},\quad L_{2,k}=E_k+A_kn^{-\alpha}+B_kd^{-\beta}+S_k\Phi_k.
\end{equation}

其中$\alpha$与$\beta$均为0.29318264389038085，领域尺度指数与注意力系数$\eta$含义不同。上游接口虽允许同比例缩放领域质量，但固定质量数据不能识别这种缩放的真实效果，故不把该接口用作治理收益。


\subsection{治理质量的条件收益}

采用第二问质量部分交付的B7表达式，以十亿为N/D单位：


\begin{equation}\label{eq:q3_5}
F_7(n,d,Q)=0.9654077057+0.5307187920n^{-0.2825}+1.6157571765(dQ)^{-0.108}.
\end{equation}

上式为显示舍入，执行使用交接文件全精度。主情景使用比值桥接，加性桥接用于判断结果对函数拼接方式的敏感性：


\begin{equation}\label{eq:q3_6}
L_k^{(r)}=L_{2,k}\left[\frac{F_7(n,d,Q)}{F_7(n,d,Q_0)}\right]^\kappa,\quad L_k^{(a)}=L_{2,k}+\kappa[F_7(n,d,Q)-F_7(n,d,Q_0)].
\end{equation}

两者在$Q=Q_0$时均回接第二问。主强度为1，正收益敏感性取0.5、1、1.5，零收益负对照取0。固定N/D/Q时，比值为正、加性增量与配比无关，两者都保持配方排序。因此，联合优化产生的非加性应解释为预算共享及资源重分配，不能声称识别了直接的质量与配方因果交互。


\subsection{成本、可行域与上下文}

三类单位Token质量成本见表~\ref{tab:q3_3}。扣除基线后只计增量，图~\ref{fig:q3_3}展示实际进入预算的增量曲线。对数形式在Q=1时有限，不能当作趋于无穷的质量屏障。


\begin{table}[htbp]\centering\setlength{\tabcolsep}{3pt}\renewcommand{\arraystretch}{1.12}
\caption{质量成本函数及题给参数}\label{tab:q3_3}
\begin{tabularx}{\textwidth}{>{\raggedright\arraybackslash}p{2.578448cm}>{\raggedright\arraybackslash}p{6.941975cm}>{\raggedright\arraybackslash}X}\toprule
形式 & $g(Q)$ & 单位与说明 \\
\midrule
指数 & $10^7\exp(6Q)$ & FLOPs/Token，随质量凸增 \\
幂函数 & $5\times10^9Q^4$ & FLOPs/Token，四次幂 \\
对数 & $2\times10^9\ln(1+10Q)$ & FLOPs/Token，上界处有限 \\
\bottomrule\end{tabularx}
\end{table}

\begin{equation}\label{eq:q3_7}
a=6+\eta L_{\mathrm{ctx}},\quad h(Q)=[g(Q)-g(Q_0)]_+,\quad\eta=2\times10^{-4}.
\end{equation}

\begin{equation}\label{eq:q3_8}
C_{\mathrm{used}}=6ND+\eta NDL_{\mathrm{ctx}}+Dh(Q)=D[aN+h(Q)].
\end{equation}

\begin{figure}[htbp]
\centering
\includegraphics[width=.95\textwidth]{figs/process_q3_quality_cost.pdf}
\caption{三种原价成本的单位Token增量h(Q)；价格量级与曲线形状同时不同}\label{fig:q3_3}
\end{figure}

设A为m行17列的可用训练锚点矩阵，配方表示为非负归一权重的凸组合。主域采用24个真实配方加训练均值，避免在完整单纯形顶点任意外推。完整训练凸包的可行证书与缩减域内的求解资格分别检查。优化模型为：


\begin{equation}\label{eq:q3_9}
\min J=\frac{1}{13}\sum_{k=1}^{13}L_k,\quad C_{\mathrm{used}}\le C,\quad p=A^\top w,\quad w\ge0,\quad\mathbf1^\top w=1.
\end{equation}

\begin{equation}\label{eq:q3_10}
10^6\le N\le10^{11},\quad10^9\le D\le10^{13},\quad Q_0\le Q\le1.
\end{equation}

所有预测必须有限且为正。数值求解内部要求各域Loss不低于$10^{-8}$；主接受规则拒绝低于$10^{-6}$的结果，正式报告点实际均严格高于该值，补充搜索采用严格正阈值。固定配比基线为512个归一训练行的算术均值，而非17域等配比。


\begin{equation}\label{eq:q3_11}
p_{\mathrm{ref}}=\frac1{512}\sum_{r=1}^{512}a_r,\quad\frac{C_{\mathrm{attn}}}{C_{\mathrm{train}}}=\frac{\eta L_{\mathrm{ctx}}}{6},\quad L_{\mathrm{ctx}}^{\mathrm{crit}}=\frac6\eta=30000.
\end{equation}

临界长度由成本解析得到，不新增为C7的第六档。固定上下文时，注意力与训练的成本比不随预算改变。表~\ref{tab:q3_4}在C为$10^{22}$FLOPs、幂成本、原25锚点联合域下比较五档上下文；单位统一为十亿。随着该项开销增加，配置的N/D均下降、预测Loss上升，但当前模型没有上下文带来的直接能力收益，不能据此建议所有任务使用最短窗口。


\begin{table}[htbp]\centering\setlength{\tabcolsep}{3pt}\renewcommand{\arraystretch}{1.12}
\caption{原25锚点联合解的上下文配置比较}\label{tab:q3_4}
\begin{tabularx}{\textwidth}{>{\raggedright\arraybackslash}p{2.351672cm}>{\raggedright\arraybackslash}p{3.476385cm}>{\raggedright\arraybackslash}p{2.556165cm}>{\raggedright\arraybackslash}p{2.556165cm}>{\raggedright\arraybackslash}p{2.044932cm}>{\raggedright\arraybackslash}X}\toprule
$L_{\mathrm{ctx}}$ & 注意力/训练比 & N（十亿） & D（十亿） & Q & J \\
\midrule
2048 & 0.06827 & 2.8157 & 448.44 & 1.0000 & 1.69176 \\
4096 & 0.13653 & 2.7163 & 439.08 & 1.0000 & 1.69523 \\
8192 & 0.27307 & 2.5438 & 422.26 & 1.0000 & 1.70169 \\
32768 & 1.09227 & 1.9152 & 353.43 & 1.0000 & 1.73153 \\
131072 & 4.36907 & 1.1338 & 245.24 & 1.0000 & 1.79558 \\
\bottomrule\end{tabularx}
\end{table}

\section{资源替代机制与数值求解}

\subsection{质量收益与数据挤占}

固定上下文、配方及参数量，若预算活跃，则D等于预算除以单位Token成本。提高质量一方面改变响应，另一方面减少可训练Token。对该预算可行方向求导：


\begin{equation}\label{eq:q3_12}
D=\frac{C}{aN+h(Q)},\quad\frac{dD}{dQ}=-\frac{Dg'(Q)}{aN+h(Q)},\quad T_Q=J_Q-\frac{J_DDg'(Q)}{aN+h(Q)}.
\end{equation}

在该方向不被其它约束阻断且响应可微时，质量下界的局部必要条件为$T_Q$非负，内点约为零，上界非正。若D已触下界、预算不活跃或存在其它活跃约束，则使用对应的可行方向，不能机械套用。端点价格归一化时将g的导数替换为归一化增量成本的导数。该判据揭示：治理是否值得取决于直接质量收益能否抵消Token挤占损失。


\subsection{规模弹性与治理成本份额}

在正目标、可微、N/D内点、预算活跃且其它约束不向N/D必要条件增加项时，拉格朗日乘子法给出两条资源驻点条件。若$J_D$非零，消去乘子得到：


\begin{equation}\label{eq:q3_13}
J_N+\lambda aD=0,\quad J_D+\lambda(aN+h)=0.
\end{equation}

\begin{equation}\label{eq:q3_14}
\varepsilon_N=-\frac{NJ_N}{J},\quad\varepsilon_D=-\frac{DJ_D}{J},\quad\frac{\varepsilon_N}{\varepsilon_D}=\frac{aN}{aN+h}=1-s_Q.
\end{equation}

无治理增量时恢复两种资源弹性相等；存在治理成本时，每增加Token还须承担处理成本，最优替代比因而出现修正项。治理份额的恒等式本身不要求预算活跃：


\begin{equation}\label{eq:q3_15}
s_Q=\frac{C_Q}{C_{\mathrm{used}}}=\frac{h}{aN+h},\quad\frac{C_Q}{C}=s_Q\frac{C_{\mathrm{used}}}{C}.
\end{equation}

\begin{equation}\label{eq:q3_16}
\left.\frac{\partial s_Q}{\partial N}\right|_{Q,L_{\mathrm{ctx}}}=-\frac{ah}{(aN+h)^2}<0\quad(h>0).
\end{equation}

因此，在质量已饱和、上下文和价格固定而N继续增大的条件下，治理份额下降完全可能，同时绝对治理成本仍可增长。该恒等关系不证明N必随预算单调增加，也不能越过N/D边界作无限预算外推。上述关系是本题成本结构下标准优化原理的条件应用，不作为新的经验标度律。


\subsection{预算参数化与锚点扩展}

主求解使用预算保持的对数坐标。先求可行质量上限，再以最小数据量求参数上限，给定N/Q后再求数据上限：


\begin{equation}\label{eq:q3_17}
Q_{\mathrm{hi}}=\max\{Q\in[Q_0,1]:D_{\min}[aN_{\min}+h(Q)]\le C\}.
\end{equation}

\begin{equation}\label{eq:q3_18}
N_{\mathrm{cap}}=\min\{N_{\max},(C/D_{\min}-h)/a\},\quad N=\exp[\ln N_{\min}+u\ln(N_{\mathrm{cap}}/N_{\min})].
\end{equation}

\begin{equation}\label{eq:q3_19}
D_{\mathrm{cap}}=\min\{D_{\max},C/(aN+h)\},\quad D=\exp[\ln D_{\min}+v\ln(D_{\mathrm{cap}}/D_{\min})],\quad u,v\in[0,1].
\end{equation}

边界舍入用下界保护处理；该变换允许不耗尽预算。锚点权重作非负归一处理，并检查原权重和误差不超过$10^{-7}$。采用SciPy实现的SLSQP求解\cite{ref3}，主实验每配置3个确定性起点，种子20260924，每起点最多450次迭代、目标容差$10^{-9}$。保留收敛且实际约束合格的最佳局部候选；不同起点不合并成全局最优证明。


主求解和算法对照的扰动合格另作如下定义：在最优候选的预算坐标u、v及Q上分别作正负扰动并截取到边界，配比可变时将锚点权重朝各顶点作凸组合扰动；步长取$10^{-4}$和$10^{-3}$。由预算映射还原试点，保留有限且满足各域Loss下限的预测；两个步长下已检查方向的最大Loss下降均不超过$10^{-6}$，才记为扰动合格。它只检查有限方向，与下面固定资源的完整512配方方向检查不同。


缩减域会遗漏完整训练域的改进方向，故另检查所有512个真实配方。固定N/D/Q，沿下式分别取两种步长；若仍有超过$10^{-6}$的改善，每轮加入至多4个配方再联合优化：


\begin{equation}\label{eq:q3_20}
p(t)=(1-t)p+ta_r,\quad r=1,\ldots,512,\quad t\in\{10^{-4},10^{-3}\}.
\end{equation}

扩展流程见图~\ref{fig:q3_4}，预设最多12轮、每轮2起点。只有重新求解的Loss改善超过$10^{-10}$才继续扫描，否则保留未解决状态。扩域与近等价补充搜索使用log10尺度N/D及显式归一预算约束，不与主参数化混称。有限方向检查通过只排除了这些已检查的小步方向，不能代替连续驻点或全局证书。


\begin{figure}[H]
\centering
\includegraphics[width=.95\textwidth]{figs/flow_q3_model.pdf}
\captionsetup{font=small}
\caption{33个重点配置的真实锚点扩展与参考解检验流程；不代表全部新推荐均已做近等价检查}\label{fig:q3_4}
\end{figure}

\section{实验设计与资源配置结果}

本节分析、代码调试和图文草拟使用OpenAI的Codex辅助，本机客户端为0.155.0-alpha.16.3；后端型号及该预发布构建的发布日期未核实。实际使用范围见文末说明。以下数值来自响应函数的条件预测优化，不代表已训练对应数量的模型。


四组机制消融见表~\ref{tab:q3_5}。规模基线只释放N/D，配比组增加p，质量组增加Q，联合组同时释放四类变量。四组共享成本、响应、边界及25锚点域；这不是独立B1标度律基线。联合可行域包含三个对照，若联合候选劣于任一对照超过$10^{-7}$，应判为求解不足，不能解释为联合策略本身较差。


\begin{table}[htbp]\centering\setlength{\tabcolsep}{3pt}\renewcommand{\arraystretch}{1.12}
\caption{四组消融与条件实验矩阵}\label{tab:q3_5}
\begin{tabularx}{\textwidth}{>{\raggedright\arraybackslash}p{2.975132cm}>{\raggedright\arraybackslash}p{7.933686cm}>{\raggedright\arraybackslash}X}\toprule
类别 & 自由变量或变化因素 & 数量及用途 \\
\midrule
规模基线 & N、D；固定平均配方与基础质量 & 与另三组共同构成主矩阵 \\
配比组 & N、D、p；固定基础质量 & 分离配比调整收益 \\
质量组 & N、D、Q；固定平均配方 & 分离治理收益及其成本 \\
联合组 & N、D、p、Q & 同域联合收益 \\
主矩阵 & 11预算、5上下文、3成本、4组 & 660项 \\
成本对照 & 统一Q0至1的端点增量价格 & 660项 \\
桥接对照 & 2形式、3强度、2基线质量 & 7920项 \\
搜索敏感性 & 49/97锚点、6起点、另种子及3种边界扰动 & 4620项 \\
零收益对照 & 质量收益强度设为0 & 660项 \\
预算加密 & 41档预算、其余主设置相同 & 2460项 \\
算法横向对照 & 5预算、2上下文、3成本、2求解器 & 另外60项 \\
\bottomrule\end{tabularx}
\end{table}

\FloatBarrier

主预算覆盖$10^{19}$至$10^{24}$的11档，含题给三个典型量级；上下文取C7全部五档。前六批共16980项登记，去重为15960个数学配置，无运行缺失。重合配置保留来源和起点信息，不把重复记录计作独立样本。


\subsection{平均收益与领域代价}

以相同预算、上下文、成本的规模基线为配对参照，定义平均改善率和联合非加性：


\begin{equation}\label{eq:q3_21}
G=100\frac{J_{\mathrm{base}}-J_{\mathrm{joint}}}{J_{\mathrm{base}}},\quad H=J_{\mathrm{quality}}+J_{\mathrm{mixture}}-J_{\mathrm{base}}-J_{\mathrm{joint}}.
\end{equation}

指数成本、4096上下文的四组轨迹见图~\ref{fig:q3_5}。全165个原联合场景的平均改善为4.1603\%至5.8086\%，说明在统一模型假设下，释放配比和治理变量具有决策价值。非加性H以百万分之一nats为容差，71组为正、47组为负、47组在容差内，并非普遍正协同。全量统计来自全部配对结果，不由单幅切片图外推。


\begin{figure}[htbp]
\centering
\includegraphics[width=.95\textwidth]{figs/result_q3_loss_ablation.pdf}
\caption{指数成本、上下文4096时11档预算的四组消融；各组均联动优化N/D}\label{fig:q3_5}
\end{figure}

均值改善伴随分配代价：全165组中162组至少一个领域退步，最多同时退步2个领域。最坏绝对退步为0.0867182 nats，出现在预算为$3\times10^{20}$、上下文2048、幂成本的dm mathematics，Loss由1.467935升至1.554654。典型9例的13域变化见图~\ref{fig:q3_6}。若实际应用要求逐域不退步，应另行加入保障约束；本文不在看到结果后改动等权目标。


\begin{figure}[htbp]
\centering
\includegraphics[width=.95\textwidth]{figs/result_q3_domain_tradeoffs.pdf}
\caption{上下文4096、三预算乘三成本的逐域Loss变化；全165组统计另按完整配对表计算}\label{fig:q3_6}
\end{figure}

\subsection{搜索域限制与有限候选配置}

完整训练配方方向检查发现，165个主联合配置均存在25锚点域外的改善方向。对上下文4096的33组执行锚点扩展，实际最多5轮，最终使用39至42个锚点，相对原联合解改善0.8383\%至2.5512\%；配比的L1变化为0.5280至0.6468。图~\ref{fig:q3_7}表明搜索域限制具有实质影响，但该扩域收益不能混入原四组机制消融。


\begin{figure}[htbp]
\centering
\includegraphics[width=.95\textwidth]{figs/result_q3_hull_expansion.pdf}
\caption{33个4096上下文配置的扩域结果相对原25锚点联合解的变化}\label{fig:q3_7}
\end{figure}

进一步对照同模型、原价、相同基础质量、原N/D边界及比值桥接强度1下的已有最终候选，33个扩域解中仍有13个不及此前搜索结果、20个更好。于是把825条主实验及搜索候选与33条扩域候选整合，逐点重算后，按每个场景的实际最小Loss选择；完全并列时按固定来源标识确定展示代表。858条记录并非858个独立物理解，也不包含所有中间迭代点。


\FloatBarrier

165组获选配置中145组来自97锚点搜索，20组来自扩域；相对原25锚点联合解进一步改善0.9583\%至2.5512\%。这表示有限候选整理避免了漏选，不构成等计算预算的算法优越性证据。表~\ref{tab:q3_6}给出上下文4096的三预算、三成本推荐，完整精度与其余上下文结果保存在电子配置表。


\begin{table}[htbp]\centering\setlength{\tabcolsep}{3pt}\renewcommand{\arraystretch}{1.12}
\caption{同合同有限候选库中的典型条件配置}\label{tab:q3_6}
\begin{tabularx}{\textwidth}{>{\raggedright\arraybackslash}p{2.487396cm}>{\raggedright\arraybackslash}p{1.989916cm}>{\raggedright\arraybackslash}p{2.288404cm}>{\raggedright\arraybackslash}p{2.387900cm}>{\raggedright\arraybackslash}p{1.691429cm}>{\raggedright\arraybackslash}p{2.089412cm}>{\raggedright\arraybackslash}X}\toprule
C（FLOPs） & 成本 & N（十亿） & D（十亿） & Q & J & 候选来源 \\
\midrule
$10^{19}$ & 指数 & 0.0678 & 21.64 & 0.6219 & 2.38753 & 扩展锚点 \\
$10^{19}$ & 幂 & 0.0678 & 21.64 & 0.6219 & 2.38753 & 扩展锚点 \\
$10^{19}$ & 对数 & 0.0678 & 21.64 & 0.6219 & 2.38753 & 扩展锚点 \\
$10^{22}$ & 指数 & 2.5115 & 491.40 & 0.9829 & 1.66994 & 扩展锚点 \\
$10^{22}$ & 幂 & 2.6199 & 452.13 & 1.0000 & 1.67192 & 扩展锚点 \\
$10^{22}$ & 对数 & 2.2175 & 626.41 & 1.0000 & 1.66062 & 97锚点 \\
$10^{24}$ & 指数 & 21.9273 & 6529.84 & 1.0000 & 1.46600 & 97锚点 \\
$10^{24}$ & 幂 & 22.0068 & 6480.02 & 1.0000 & 1.46612 & 97锚点 \\
$10^{24}$ & 对数 & 21.5744 & 6758.48 & 1.0000 & 1.46547 & 97锚点 \\
\bottomrule\end{tabularx}
\end{table}

对于表~\ref{tab:q3_6}，低预算三种成本均保持$Q=Q_0$，因为此时没有治理增量，成本形式不再区分配置；中高预算的质量提升须结合各自代价解释。表~\ref{tab:q3_7}列出同一有限库、幂成本、上下文4096的完整17维配方，三个预算列分别与表~\ref{tab:q3_6}的幂成本来源对应。


\begin{table}[H]\centering\setlength{\tabcolsep}{3pt}\renewcommand{\arraystretch}{1.12}
\caption{幂成本条件下典型有限候选的完整领域配方（\%）}\label{tab:q3_7}
\begin{tabularx}{\textwidth}{>{\raggedright\arraybackslash}p{5.908112cm}>{\raggedright\arraybackslash}p{3.249462cm}>{\raggedright\arraybackslash}p{3.249462cm}>{\raggedright\arraybackslash}X}\toprule
领域 & $10^{19}$ & $10^{22}$ & $10^{24}$ \\
\midrule
arxiv & 3.45 & 2.62 & 1.79 \\
freelaw & 7.67 & 7.44 & 5.25 \\
nih exporter & 0.02 & 0.01 & 4.88 \\
pubmed central & 2.93 & 1.46 & 3.63 \\
wikipedia en & 6.84 & 7.74 & 19.69 \\
dm mathematics & 4.73 & 5.06 & 0.84 \\
github & 11.40 & 12.92 & 12.12 \\
philpapers & <0.005 & <0.005 & <0.005 \\
stackexchange & 14.07 & 13.10 & 3.83 \\
enron emails & 0 & <0.005 & 0.08 \\
gutenberg pg 19 & 6.55 & 9.03 & 15.90 \\
pile cc & 17.87 & 28.14 & 11.49 \\
ubuntu irc & 5.10 & 5.13 & 14.81 \\
europarl & <0.005 & <0.005 & 0.09 \\
hackernews & 6.16 & 4.63 & 5.22 \\
pubmed abstracts & 9.81 & 0.05 & <0.005 \\
uspto backgrounds & 3.41 & 2.67 & 0.36 \\
\bottomrule\end{tabularx}
\end{table}

完整配方在表~\ref{tab:q3_7}中以百分比表示，独立四舍五入到两位小数；小于0.005\%的严格正值显示为<0.005，真零显示0，显示值之和不保证恰为100\%。源结果未据此重新归一化。新推荐中20个确切扩域来源保有完整512配方两步长记录，另145个仅有其缩减域扰动记录；近等价审查仅覆盖5个确切扩域参考，其余160个未检查。推荐表不据此宣称配方唯一或全域最优。


\section{成本敏感性与结构性转移}

\subsection{治理份额与价格尺度}

实际成本份额见图~\ref{fig:q3_8}，其范围为原25锚点、上下文4096的9个典型配置。质量提高、治理绝对成本增加、治理份额上升是不同命题。另以扩域的幂成本案例说明：预算由$10^{22}$增至$10^{24}$时，Q均为1，治理份额由19.2241\%降至2.6935\%，但绝对治理成本由$1.9224\times10^{21}$增至$2.6935\times10^{22}$FLOPs，即增至原来的约14.01倍。这组扩域数值不作为图中原域数值。


\begin{figure}[htbp]
\centering
\includegraphics[width=.95\textwidth]{figs/result_q3_cost_allocation.pdf}
\caption{原25锚点、上下文4096、三预算乘三成本的训练、注意力及治理份额}\label{fig:q3_8}
\end{figure}

为检验资源替代条件，采用两种相对差分步长计算归一化残差：


\begin{equation}\label{eq:q3_22}
r_{ND}=\frac{NJ_N-aNDJ_D/(aN+h)}{\max(|J|,1)},\quad\Delta\in\{10^{-5},10^{-4}\}.
\end{equation}

165个原联合解共330次检查的最大绝对残差为$2.04434\times10^{-6}$；图~\ref{fig:q3_9}仅展示步长为$10^{-5}$的165点。小残差支持推导与模型内求解一致，不能证明响应能准确预测真实训练。质量边界的预算轨迹见图~\ref{fig:q3_10}，连续颜色本身不是结构转移判定。


\begin{figure}[htbp]
\centering
\includegraphics[width=.95\textwidth]{figs/process_q3_resource_balance.pdf}
\caption{原联合解的资源均衡检查；图中为相对差分步长10的负5次方的165点}\label{fig:q3_9}
\end{figure}

\begin{figure}[htbp]
\centering
\includegraphics[width=.95\textwidth]{figs/result_q3_quality_regimes.pdf}
\caption{原联合域在五档上下文、三成本及11预算下的连续质量值；色阶不代表稳定相变标签}\label{fig:q3_10}
\end{figure}

题给三成本同时改变价格和形状。为了分离二者，在同一Q0下把Q=1的增量对齐到指数成本：


\begin{equation}\label{eq:q3_23}
\lambda_f=\frac{g_{\exp}(1)-g_{\exp}(Q_0)}{g_f(1)-g_f(Q_0)},\quad h_f^{\mathrm{norm}}(Q)=\lambda_fh_f(Q).
\end{equation}

归一化前后质量轨迹见图~\ref{fig:q3_11}。对55个预算/上下文组，要求两种设定中排名第一与第二的Loss差均超过百万分之一，仍有19组最优成本形式改变。因此，原价下哪种成本给出更低Loss，不能全部归因于曲线形状。


\begin{figure}[htbp]
\centering
\includegraphics[width=.95\textwidth]{figs/result_q3_quality_cost_normalized.pdf}
\caption{上下文4096时原价与端点价格归一化后的Q轨迹；19/55排名变化另由完整比较表计算}\label{fig:q3_11}
\end{figure}

\subsection{结构状态、机制与持续性}

把结构性转移操作化为质量边界状态、有效配方支持或N/D活跃边界的变化，并要求这种变化通过机制、持续性及情景检验。质量距下界不超过0.001记为下界状态，距1不超过0.001记为上界，否则为内点；配比状态定义为：


\begin{equation}\label{eq:q3_24}
S_\tau(p)=\{i:p_i>\tau\},\quad\tau=0.01,\quad d_J(S_0,S_1)=1-\frac{|S_0\cap S_1|}{|S_0\cup S_1|}.
\end{equation}

并集为空时距离记0。配方候选变化须同时满足$d_J\ge0.2$和$\|p_1-p_0\|_1\ge0.02$。两端必须数值合格；质量机制检查预算方向量，方向容差为$10^{-4}$；配方还要求已接受的近等价起点没有不同支持。后续至少两个预算点维持新状态，再要求同一粗预算区间内的加密数据严格收缩变化区间。任一成本份额变化达到0.15只作描述，不单独判定机制转移。


按同成本、上下文和优化模式，将同机制、区间严格重叠且通过机制与持续检验的轨迹数除以全部登记轨迹数，作为情景支持比例；需覆盖桥接、强度、种子、边界及锚点扰动。主比例门槛为70\%，另看60\%和80\%。近等价起点的绝对Loss容差默认$10^{-4}\max(|J|,1)$，其中$10^{-4}$是无量纲缩放系数，并检查阈值敏感性。比例描述注册假设下的一致性，不是统计置信水平。


结构识别结果见表~\ref{tab:q3_8}。固定配方、指数成本、最长上下文仅有一个条件质量启动事件，其粗细区间表示同一事件。联合策略没有通过70\%规则的稳定质量或有效支持事件，因此不能用少量典型配置的差异宣称普遍联合相变。


\begin{table}[H]\centering\setlength{\tabcolsep}{3pt}\renewcommand{\arraystretch}{1.12}
\caption{结构转移识别与阈值敏感性}\label{tab:q3_8}
\begin{tabularx}{\textwidth}{>{\raggedright\arraybackslash}p{3.570159cm}>{\raggedright\arraybackslash}p{7.636173cm}>{\raggedright\arraybackslash}X}\toprule
项目 & 实际结果 & 判断 \\
\midrule
固定配方质量启动 & 指数成本，上下文131072 & 一个下界至内点事件 \\
预算区间（FLOPs） & 粗：$[3\times10^{19},10^{20}]$；细：$[7.49894\times10^{19},10^{20}]$ & 细区间属于同一事件 \\
情景支持比例 & 16/21，即76.19\% & 70\%通过，80\%不通过 \\
质量方向容差 & $10^{-4}$通过，$10^{-5}$不通过 & 事件依赖数值门槛 \\
联合配置 & 未检出通过70\%规则的稳定事件 & 不宣称普遍结构相变 \\
光滑配方负对照 & 所有占比严格为正仍可跨1\%阈值 & 有效支持改变不等于零支持突变 \\
\bottomrule\end{tabularx}
\end{table}

\subsection{近等价配方与识别边界}

原域与扩域各选上下文4096、三个典型预算乘三成本的9个参考，分别保留自身支持域和目标基准。对每个参考定义近等价集合，选择最靠近1\%阈值的三个领域，分别最小化、最大化占比：


\begin{equation}\label{eq:q3_25}
\mathcal E_\epsilon=\{x\in\mathcal F_{\mathrm{ref}}:J(x)\le\widehat J+\epsilon\},\quad\epsilon=\rho\max(|\widehat J|,1),\quad\rho\in\{10^{-5},10^{-4},10^{-3}\}.
\end{equation}

每任务至多2起点，每起点最多400次SLSQP迭代，目标容差$10^{-10}$。原域按目标值、科学内容标识和原起点序号排序选取已接受端点，保留其重数；扩域采用最终扩域权重与均匀扩域权重。两批各自重新选取近阈值领域，搜索预算相同不表示参考目标或可行域相同。


原域和扩域各执行162个任务、324次起点求解，分别保留7次和4次未收敛。原域9个任务发现跨阈值反例，涉及6个典型配置；扩域1个任务发现反例，涉及1个配置。相同参考、领域、支持域和阈值下，小容差见证可继承到大容差，已知反例集合分别为15/81和1/81。两次搜索的参考和领域不同，不能用9比1证明扩域提高了稳定性；没有找到反例也不证明唯一。


扩域的反例位于$C=10^{24}$、上下文4096、对数成本：允许约0.1\%的平均Loss增加，pubmed abstracts占比可由0.0410924\%升至3.5257574\%，质量仍近上界。它是近等价而非更优解。相对后来有限库选出的更低Loss参考，该见证增幅约0.22592\%，已超过0.1\%上限，不能转移为新推荐的同容差反例。这个例子说明结构判定必须绑定参考目标、可行域和容差。


\section{数值检验、不确定性与模型评价}

\subsection{有限情景下的稳健选择}

在共同基础质量与可比成本下，对两种桥接、三种正收益强度共六情景互评共同可行的有限候选。令X为经过复核的该组有限集合，定义各情景相对后悔和推荐配置：


\begin{equation}\label{eq:q3_26}
J_s^{\mathrm{best}}=\min_{x\in X}J_s(x),\quad R_s(x)=\frac{J_s(x)-J_s^{\mathrm{best}}}{|J_s^{\mathrm{best}}|}.
\end{equation}

\begin{equation}\label{eq:q3_27}
x_{\mathrm{rob}}\in\mathop{\arg\min}_{x\in X}\max_sR_s(x),\quad v_{\mathrm{rob}}=\min_{x\in X}\max_sR_s(x).
\end{equation}

所有场景Loss为正；实际值完全并列时按固定候选标识选择展示代表，近并列标志不改变真正取最小的规则。全165组的最小最大相对后悔值最高1.5465\%，图~\ref{fig:q3_12}展示其中4096上下文的33组。把相同推荐另置于零收益压力情景，最高后悔为2.5946\%；零收益基准使用正收益候选与零收益候选的并集，但没有参与原推荐选择。该结论只针对已互评的候选及情景，不转移到后续有限库的新推荐，也不是连续域稳健保证。


\begin{figure}[htbp]
\centering
\includegraphics[width=.95\textwidth]{figs/result_q3_finite_regret.pdf}
\caption{上下文4096的33组有限候选六情景后悔值；全165组最大值由完整互评计算}\label{fig:q3_12}
\end{figure}

零收益负对照中，330个自由Q配置均返回基线质量，最大偏差约$9.50\times10^{-14}$。理论上，当$\kappa=0$，把任一可行点的Q替换为$Q_0$保持目标且成本不增；最优解若存在，至少存在基线质量的最优解。只有当节约成本还能换取可行的严格目标改善，例如D可增且$J_D<0$时，才可进一步排除较高Q的最优解，不能无条件宣称策略唯一。


\subsection{求解器与复算一致性}

SLSQP与COBYQA的30对比较使用相同25锚点域、相同3起点，每起点最多3000次内部响应核调用及4000次迭代；这与主矩阵450次迭代设置分开。后处理调用单独记录，未给单方追加额度。结果见图~\ref{fig:q3_13}：60项均收敛可行，但扰动合格为59项，SLSQP为30/30，COBYQA为29/30。


\begin{figure}[htbp]
\centering
\includegraphics[width=.88\textwidth]{figs/process_q3_solver_agreement.pdf}
\caption{30对同起点、同支持域、同调用上限的求解器比较；扰动不合格例保留}\label{fig:q3_13}
\end{figure}

\FloatBarrier

两算法最大Loss差为0.000168474，来自预算$10^{22}$、上下文2048、幂成本的一例；该COBYQA点可行扰动仍能改善$1.64197\times10^{-6}$，超过登记阈值。对每配置先合计三个起点的内部响应调用，再在30个配置上取中位数，两方法分别为2715.5和4734；这不是每起点调用数或运行时间，只说明这些问题及实现中的调用差异。候选的成本、配比和13域Loss另行重算；最终有限库最大相对预算超额为$9.43362\times10^{-11}$，在固定的$10^{-10}$接受容差内，不能表述成严格零超额。


\subsection{适用边界与后续实证}

本模型的优势是统一三项算力成本，给出可检验的资源替代判据，显式限制配比支持域，并把平均收益、领域损害与策略可识别性一起报告。主要限制来自证据而非优化器：B7半合成质量标尺与本题治理质量尚未独立校准，第二问缺少完整拟合程序、训练/留出ID、参数协方差及B1平行基线参数；规模范围属于条件外推。因此，不报告统计显著性、预测置信区间或实际训练改善率。


可执行的后续实证应在同一训练流水线上，对典型预算各选择基线、单项及联合配置，用相同评测集和随机种子组复测13域Loss，同时记录质量处理的真实算力。先校准治理强度和单位Token成本，再判断配对收益及领域退步是否保留；若校准显示收益很弱，应比较基础质量方案与对应弱收益情景下的可行配置，再确定选择。上述校准与训练尚未实施，不把已完成的压力检验当作另一个已交付策略。


\clearpage

\section{结论}

在前两问响应与明确治理假设下，建立了N、D、17域配比和质量的联合配置模型，给出三个预算量级及五档上下文的配置结果。注意力与训练成本相当的解析上下文长度为30000；当前模型仅刻画其成本影响。原同域联合配置的平均Loss相对内部规模基线改善4.16\%至5.81\%，但多数场景伴随个别领域损失，均值目标不能替代领域保障。


质量治理同时带来直接收益与Token挤占，资源弹性比中的治理份额修正解释了这一替代关系。质量饱和后，其算力份额可以下降而绝对投入继续增长。成本比较中19/55组排名受端点价格归一化影响，表明价格与形状必须分开。按既定持续性、加密及70\%情景支持规则，未发现稳健的联合结构转移；固定配方仅存在一个对阈值敏感的质量启动事件。最终建议以有限候选配置作为条件决策起点，保留其外推与近等价边界，经实际训练验证后再作部署选择。

\begin{thebibliography}{99}
\bibitem{ref1} Hoffmann J, Borgeaud S, Mensch A, et al. An empirical analysis of compute-optimal large language model training[C]. Advances in Neural Information Processing Systems, 35, 2022. DOI: 10.52202/068431-2176.
\bibitem{ref2} Liu Q, Zheng X, Muennighoff N, et al. RegMix: Data Mixture as Regression for Language Model Pre-training[C]. International Conference on Learning Representations, 2025.
\bibitem{ref3} Virtanen P, Gommers R, Oliphant T E, et al. SciPy 1.0: Fundamental Algorithms for Scientific Computing in Python[J]. Nature Methods, 2020, 17: 261-272. DOI: 10.1038/s41592-019-0686-2.
\end{thebibliography}
\section*{AI工具使用说明}
本研究参考分稿使用OpenAI开发的Codex辅助方案衔接讨论、数学推导整理、程序生成与调试、实验执行结果整理、图表制作和文字草拟。本机可验证的客户端版本为0.155.0-alpha.16.3；客户端版本不同于后端模型型号，当前具体后端型号与该预发布构建的发布日期未核实，未借用稳定版本发布日期代填。

数学论断以文中推导、交接数据和已列文献为依据，数值来自实际保存的条件优化结果。本文没有实际训练所列全部配置，也不声称参赛者已经独立完成、理解并人工审核全文。正式使用前仍须由使用者核对、理解并以自己的语言表述，补齐工具信息，完成全题整合与竞赛要求的最终审核；本分稿不等同于已获正式提交合规认定的作品。
\end{document}
